\chapterhead{40}{ORR-1}{Introduction to Operational Risk and Resilience}{4}{0}

\lohead{40}{a}{Describe an operational risk management framework and assess the
types of risks that can fall within the scope of such a framework.}

\begin{ordefbox}{Definition}
Operational risk has been defined by the Basel Committee on Banking Supervision
(BCBS) as ``the risk or loss resulting from inadequate or failed internal
processes, people, systems, and external events.''
\end{ordefbox}

\noindent{\sffamily\bfseries\fontsize{9.6}{12}\selectfont\color{navy}
The four causes of operational risk:}\ {\fontsize{10}{13}\selectfont
inadequate or failed internal processes, people, systems, and external events.}

\vspace{10pt}
\begin{center}
\begin{tikzpicture}[font=\fontsize{8.4}{10.5}\selectfont,
  band/.style={text width=4.0cm,align=left,rounded corners=3pt,inner sep=6pt,
               text=white,font=\sffamily\bfseries\fontsize{9}{11}\selectfont},
  desc/.style={text width=10.3cm,align=left,inner sep=6pt,text=navy}]

\node[band,fill=navy,anchor=north west]  (s1) at (0,0)      {Risk Identification};
\node[band,fill=teal,anchor=north west]  (s2) at (0,-1.55)  {Risk Assessment};
\node[band,fill=forest,anchor=north west](s3) at (0,-3.10)  {Risk Mitigation};
\node[band,fill=olive,anchor=north west] (s4) at (0,-4.65)  {Risk Monitoring};

\node[desc,anchor=west] at ($(s1.east)+(0.3,0)$)
  {Key risks to your objectives and strategy.};
\node[desc,anchor=west] at ($(s2.east)+(0.3,0)$)
  {Evaluate the risks, in likelihood and severity, to prioritise the response.};
\node[desc,anchor=west] at ($(s3.east)+(0.3,0)$)
  {What to do to prevent incidents and to reduce their impact.};
\node[desc,anchor=west] at ($(s4.east)+(0.3,0)$)
  {Check if everything works as intended, learn the lessons if not, and monitor
   the evolution of risks.};

\foreach \a/\b in {s1/s2,s2/s3,s3/s4}
  \draw[-{Stealth[length=5pt]},navy!35,line width=1.6pt]
    ($(\a.south west)+(0.55,0)$) -- ($(\b.north west)+(0.55,0)$);

\draw[-{Stealth[length=4pt]},rust,line width=1pt,dashed]
  ($(s4.south west)+(0.55,0)$) .. controls +(0,-0.55) and +(0,-0.55) ..
  ($(s4.south west)+(-0.85,0)$) -- ($(s1.north west)+(-0.85,0)$)
  .. controls +(0,0.5) and +(0,0.5) .. ($(s1.north west)+(0.55,0)$);
\node[text=rust,font=\sffamily\fontsize{7.4}{9}\selectfont,rotate=90,anchor=south]
  at ($(s2.south west)+(-1.0,-0.3)$) {continuous cycle};
\end{tikzpicture}
\end{center}
\orfigcap{The operational risk management cycle.}

\orsub{Types of risks within the ORM framework}

\begin{lettered}
\item \textbf{Legal / compliance:} failing to follow rules, internal or external,
leads to fines and reputational damage.
\item \textbf{Reputational:} operational incidents or strategic choices can harm
a firm's public image.
\item \textbf{Strategic:} poor decisions, or weak execution of good plans, can
cause financial losses.
\end{lettered}

% ---------------------------------------------------------------------
\lohead{40}{b}{Describe the seven Basel II event risk categories and identify
examples of operational risk events in each category.}

\noindent\begin{minipage}{\textwidth}
\renewcommand{\arraystretch}{1.35}
\begin{tabularx}{\textwidth}{@{}L{0.30\textwidth}YL{0.14\textwidth}L{0.15\textwidth}@{}}
\rowcolor{navy} \thd{Category} & \thd{Examples} & \thd{Frequency} & \thd{Loss severity}\\
{\tabfont \textbf{1. Internal fraud (IF)}} & {\tabfont Employee defalcation, rogue trading} & \chiplow & \chiplow\\
\rowcolor{paleblue}
{\tabfont \textbf{2. External fraud (EF)}} & {\tabfont Credit card fraud, losses from hacking} & \chiphigh & \chiplow\\
{\tabfont \textbf{3. Employment practices and workplace safety (EPWS)}} & {\tabfont Employee termination, discrimination} & \chipmod & \chiplow\\
\rowcolor{paleblue}
{\tabfont \textbf{4. Clients, products, and business practices (CPBP)}} & {\tabfont Client errors, regulatory fines} & \chiphigh & \chipvhigh\\
{\tabfont \textbf{5. Damage to physical assets (DPA)}} & {\tabfont Weather-related events, negligence} & \chiplow & \chiplow\\
\rowcolor{paleblue}
{\tabfont \textbf{6. Business disruption and system failures (BDSF)}} & {\tabfont IT problems, service interruptions} & \chiplow & \chiplow\\
{\tabfont \textbf{7. Execution, delivery, and process management (EDPM)}} & {\tabfont Clerical errors, insufficient documentation} & \chiphigh & \chiphigh\\
\end{tabularx}
\end{minipage}
\orfigcap{The seven Basel II Level 1 event-type categories.}

% ---------------------------------------------------------------------
\lohead{40}{c}{Explain characteristics of operational risk exposures and
operational loss events, and challenges that can arise in managing operational
risk due to these characteristics.}

\orsub{Operational risk characteristics}

\begin{orkeybox}
{\sffamily\bfseries\fontsize{9.4}{11}\selectfont\color{navy}Heterogeneous}\par
\vspace{2pt}
Operational risks vary widely, from minor errors with no financial loss to
significant events causing substantial damage. This diversity requires careful
categorisation and analysis.\par
\vspace{3pt}
{\fontsize{9}{11}\selectfont\color{rust}\itshape E.g.\ minor errors in data
entry, cyberattacks, natural disasters.}
\end{orkeybox}

\begin{orkeybox}
{\sffamily\bfseries\fontsize{9.4}{11}\selectfont\color{navy}Idiosyncratic}\par
\vspace{2pt}
These risks are diffuse and cannot be centralised, often managed on an
individual level by employees through robust controls and procedures. Despite
those efforts, some level of risk always remains due to their unique nature.\par
\vspace{3pt}
{\fontsize{9}{11}\selectfont\color{rust}\itshape E.g.\ company-specific risks.}
\end{orkeybox}

\begin{orkeybox}
{\sffamily\bfseries\fontsize{9.4}{11}\selectfont\color{navy}Heavy tailed}\par
\vspace{2pt}
Characterised by many minor losses and a few major ones, operational risks
exhibit significant asymmetry in loss distribution. The challenge lies in
managing the potential for large, infrequent losses.\par
\vspace{3pt}
{\fontsize{9}{11}\selectfont\color{rust}\itshape E.g.\ minor service fees versus
rogue trading incidents.}
\end{orkeybox}

\begin{orkeybox}
{\sffamily\bfseries\fontsize{9.4}{11}\selectfont\color{navy}Interconnected}\par
\vspace{2pt}
Operational risks are often correlated due to common causes like human error or
external events, and can interact with financial risks, complicating
quantification and management.\par
\vspace{3pt}
{\fontsize{9}{11}\selectfont\color{rust}\itshape E.g.\ control weaknesses;
trading errors (operational) can lead to market losses (financial).}
\end{orkeybox}

\begin{orkeybox}
{\sffamily\bfseries\fontsize{9.4}{11}\selectfont\color{navy}Dynamic}\par
\vspace{2pt}
As business practices and the external environment evolve, so do operational
risks, necessitating a flexible and sometimes reactive risk management approach
to address unexpected changes.\par
\vspace{3pt}
{\fontsize{9}{11}\selectfont\color{rust}\itshape E.g.\ regulatory fines or a
shift in regulation.}
\end{orkeybox}

% ---------------------------------------------------------------------
\lohead{40}{d}{Describe operational resilience, identify the elements of an
operational resilience framework, and summarize regulatory expectations for
operational resilience.}

\begin{ordefbox}{Operational resilience}
Operational resilience refers to how firms and industries deal with business
disruptions. It includes activities such as anticipating, reacting to, and
recovering from such disruptions.
\end{ordefbox}

\orsub{Components of operational resilience}

\begin{checks}
\item \textbf{Business continuity:} minimising disruptions to business processes.
\item \textbf{Key services:} identifying and ensuring critical services continue
with minimal interruption.
\item \textbf{Impact tolerance:} defining acceptable disruption times for key
services.
\item \textbf{Disruption processes:} responding to disruptions, maintaining
stakeholder confidence, and communicating effectively.
\item \textbf{Feedback:} learning from past incidents to prevent future
occurrences.
\end{checks}

\orsub{Regulatory expectations}

Focus has shifted from solely preventing cyber incidents to managing them
effectively, and with it the importance of operational resilience in maintaining
IT service continuity after cyberattacks.

\vspace{4pt}
\noindent\begin{minipage}{\textwidth}
\renewcommand{\arraystretch}{1.4}
\begin{tabularx}{\textwidth}{@{}L{0.22\textwidth}Y@{}}
\rowcolor{navy} \thd{Jurisdiction} & \thd{Emphasis}\\
{\tabfont \textbf{1. U.K.}} & {\tabfont Focuses on IT service continuity and adjustments for work-from-home environments.}\\
\rowcolor{paleblue}
{\tabfont \textbf{2. U.S.}} & {\tabfont Emphasises operational resilience as the main outcome of an effective ORM framework.}\\
{\tabfont \textbf{3. BCBS}} & {\tabfont Issued seven principles for operational resilience, including governance, risk management, incident management, and ICT.}\\
\rowcolor{paleblue}
{\tabfont \textbf{4. EU (DORA Act, proposed)}} & {\tabfont Aims to promote digital finance while managing associated risks; focuses on IT requirements for financial institutions.}\\
{\tabfont \textbf{5. Singapore (MAS \& ABS, 2021)}} & {\tabfont Addresses operational resilience in remote work settings; emphasises controls and employee education.}\\
\end{tabularx}
\end{minipage}
\orfigcap{Regulatory emphasis on operational resilience, by jurisdiction.}
